We show that gravity can remove the linear magnetic instability of a holographic
plasma with an $SU(2)$ global symmetry. In the probe limit, charged vector
bosons become tachyonic in their lowest Landau level above a critical field
and condense into a vortex lattice. In Einstein--$SU(2)$ Yang--Mills theory
in $\AdS_5$ we include the backreaction of the field exactly, expanding only
in the condensate around the D'Hoker--Kraus magnetic brane. The critical
field diverges as the coupling $\alpha = \kappa^2/\hat g^2$ approaches
$\astar = 2/3$, above which the normal phase has no linear charged
instability at any field or temperature: we show this at zero temperature
and establish it numerically at finite temperature. At zero temperature the
field cancels from the mass of the lowest Landau level in the brane's
$\mathrm{AdS}_3 \times \mathbb{R}^2$ throat, and the instability is a
violation of the throat's Breitenlohner--Freedman bound. The onset
temperature vanishes at $\astar$ with an essential singularity of
Berezinskii--Kosterlitz--Thouless type. For two-derivative Einstein--Yang--Mills duals in $d = 4$, 5 and 6
dimensions the normal phase has a charged instability if and only if
$C_J/C_T < 8/(d^2(d+1))$, one tenth in $d = 4$. Above $\alpha_L = 0.372$
the gravitational attraction between vortices makes the uniform lattice
unstable and the transition first order, so $B_c$ is only a spinodal;
whether a condensed phase survives above $\astar$ is left open.
